% !TeX root = ../main.tex
\clearpage

\section{Results}

All results below come from continuous multi-electrode recordings of two adult captive \textit{Electrophorus voltai} (one male, one female) in the Museum f{\"u}r Naturkunde Berlin tank.
Most sessions used a 16-channel electrode line (25\,cm spacing, 3.75\,m total length) and span November 2023 to March 2026 (165 recording days; Table~\ref{tab:recording_inventory}).
Figure~\ref{fig:results_intro_excerpt} shows a short excerpt of one electrode channel with both low- and high-amplitude EODs in the same window.

\begin{figure}[H]
  \centering
  \includegraphics[width=0.92\textwidth]{detection/results_intro_recording_excerpt.png}
  \caption{One-second excerpt from a Berlin tank recording (single electrode channel).
  Low- and high-amplitude discharges appear in the same window.
  Both enter the downstream counts and shape classes; there is no amplitude or volley filter.}
  \label{fig:results_intro_excerpt}
\end{figure}

\subsection{Special pulse shapes}

Waveform shape can reveal signalling diversity that discharge rate and amplitude alone do not capture.
We therefore examined the shapes of all detected EODs to test whether the long-term Berlin recordings contain structured waveform variation beyond a single stereotyped pulse form.
We detected $n = 13{,}465{,}438$ EODs in the processed Berlin dataset.
Throughout this thesis, ``all pulses'', ``all EODs'', and related totals mean every accepted detection, including low-voltage search pulses and high-voltage discharges (shocks/volleys); events are not split or excluded by amplitude or behavioural context.
The PCA plus random forest (RF) pipeline assigns each accepted EOD to one of three shape labels---normal, wide, or double (Figure~\ref{fig:pulse_shape_distribution}). 
Median prototypes were calculated from $20{,}000$ random samples of each label to summarise the morphological extremes of the labelled set (Figure~\ref{fig:prototypes}). 
Normal pulses make up the biggest share with 64.0\%  and a total of $n = 8{,}623{,}379$ pulses. 
The normal pulse is a single, roughly symmetric peak with median half width of $1.58$\,ms. 
Wide pulses make up about a third of the dataset (33.5\%, $n = 4{,}514{,}438$). 
A typical wide pulse is a single peak with a slower decay and median half width of $2.19$\,ms. 
The decaying part of a wide pulse is often characterized by a little "bump" or shoulder, which is not present in the normal pulse. 
Therefore wide pulses are also less symmetrical compared to normal pulses.
Double pulses are quite rare and only make up 2.4\% ($n = 327{,}621$) of all recorded pulses. 
Average double pulses have two peaks of similar height separated by a shallow trough.
In the labelling process however, the goal was to find as many double pulses as possible to improve classifier performance. 
To this end, pulses that did not have the typical morphology of two clearly elevated peaks but a more or less pronounced shoulder were also labelled as double pulses.
Peak to peak separation on the median double is $1.00$\,ms. 
Median half width is $2.40$\,ms when the waveform is treated as one event.
This makes doubles the longest of the three classes, in median almost 1 ms longer than normal pulses.  
But the half width is not a sufficient label by itself because the classes overlap in width.

\begin{figure}[H]
  \centering
  \includegraphics[width=0.72\textwidth]{pulse_shapes/pulse_shape_distribution.png}
  \caption{Number of EODs in each waveform cluster.
  Counts and percentages are displayed on the bars.
  The total includes all accepted detections regardless of amplitude (low- and high-voltage discharges).
  Most discharges are normal, about a third are wide, and doubles are uncommon.}
  \label{fig:pulse_shape_distribution}
\end{figure}


\begin{figure}[H]
  \centering
  \includegraphics[width=0.82\textwidth]{pulse_shapes/prototype_pulses_stacked.png}
  \caption{Median EOD waveforms for the three classes.
  The thick coloured trace is the class median obtained from an $n = 20{,}000$ subsample.
  The thin grey pulses are 100 example pulses per class.
  The stars mark where the waveforms were aligned (peak for normal and wide, trough between peaks for double).
  Normals are short and single peaked, wides stretch out with a slower falling edge, and doubles carry two peaks about 1\,ms apart.}
  \label{fig:prototypes}
\end{figure}

Half width distributions show graded shifts among the three labels rather than cleanly separated modes (Figure~\ref{fig:half_width}). 
Normal pulses aggregate near $\sim$1.5 to 1.7\,ms (panel median $1.56$\,ms). 
Wide pulses shift mass a little longer (median $1.94$\,ms, mode near $\sim$2\,ms). 
Doubles show least variability in half width distribution.
The double pulses form a tight peak near 2.3 to 2.5\,ms (median $2.38$\,ms). 
The $0.7$\,ms half-width threshold was applied to drop near-zero and otherwise spurious widths that mostly reflect noise or misclassified events rather than genuine EOD morphology.
Importantly, the left-panel $n$ values are not the full class counts from Figure~\ref{fig:pulse_shape_distribution}: they count only pulses with a measurable half width of at least $0.7$\,ms.
Most double waveforms do not yield a finite half-width estimate (only $n = 24{,}031$ of $327{,}621$ doubles), so their legend $n$ is much smaller than the class total.
Normal and wide share substantial support around 1.5 to 2\,ms, and the labelled distributions grade into one another, so width alone does not delineate the classes. 
The pooled half width histogram without class labels is dominated by a mode near $\sim$1.6\,ms with a long shoulder toward 2 to 2.5\,ms (Figure~\ref{fig:half_width}, right), 
consistent with overlapping, continuously varying widths rather than three discrete peaks.
In the pooled histogram the peak of double pulses at about 2.5 ms is not as pronounced as in the individual class histograms. 
This is because the normal and wide pulses occur much more frequently compared to the doubles, which causes their distributions to dominate the overall shape of the pooled histogram. 
The smaller number of double pulses results in a less distinct peak in the combined distribution.
Although we can't observe a clear bimodal distribution in the pooled histogram, the presence of the shoulder indicates that there is a subset of pulses with longer half widths. 

\begin{figure}[H]
  \centering
  \includegraphics[width=\textwidth]{pulse_shapes/half_width_distributions_combined.png}
  \caption{Half-width distributions for the three waveform classes.
  Left: class-wise densities (normal blue, wide green, double red); legend $n$ is pulses with measurable half width $\geq 0.7$\,ms (not the full RF class totals).
  Right: all classes pooled as a histogram with KDE overlay.
  Pulses shorter than $0.7$\,ms are excluded to suppress noise / misclassification artefacts; the $x$-axis stops at $3.5$\,ms.
  Normal and wide share support around 1.5 to 2\,ms, while doubles form a narrow ridge near 2.4\,ms.
  The pooled mode sits near 1.6\,ms with a shoulder toward 2 to 2.5\,ms.}
  \label{fig:half_width}
\end{figure}


A PCA of the classified waveforms shows a similar picture of overlap rather than three cleanly separated clusters (Figure~\ref{fig:pca_space}). 
PC1 and PC2 together explain 84.8\% of variance in the plotted sample (63.3\% on PC1 and 21.5\% on PC2). 
Doubles form a tightly packed cloud, while normals spread more widely along PC2. 
Wide pulses sit between them and partially overlap the double region; their scatter is closer to that of normals than of doubles. 
In the PC3/PC4 views, doubles themselves split into two clumps along PC1. 
Taken together with the half-width distributions, the PCA cloud is consistent with a morphological continuum in which labelled classes occupy overlapping regions rather than discrete islands.

\begin{figure}[H]
  \centering
  \includegraphics[width=0.95\textwidth]{pulse_shapes/labeled_pulses_pca_space.png}
  \caption{RF labelled pulses in PCA space, 5000 examples per class. 
  Blue points represent normal pulses, green wide pulses, and red double pulses. 
  The four panels are different PC pairings. 
  Doubles form two dense clouds; normals show higher variability, especially along PC2. 
  Wides lie between normal and double regions and partially overlap the double cloud.}
  \label{fig:pca_space}
\end{figure}

A long standing question since we found evidence of double pulses was how they are produced physiologically.
Aligning two of the median normal pulses with the peaks of the median double pulse sparked the hypothesis that a double peaked pulse could be the result of two normal pulses fired in quick succession.  
  To test whether a double pulse is simply two normal pulses fired within 1 ms, the median double was modelled as the sum of two median normal templates (Section~\ref{sec:waveform_classification}).
When two median normal pulses are aligned to the two peaks of the median double, they already resemble the double envelope with surprising precision (Figure~\ref{fig:normal_aligned_double}).
Several nested models then fit amplitudes $a_1$ and $a_2$, peak times $t_1$ and $t_2$ (hence $\Delta t$), and component width scales to the median double (Table~\ref{tab:double_pulse_models}; Figure~\ref{fig:double_model_hierarchy}).
Constraints on $\Delta t$ and the width scales were loosened in each subsequent model; amplitudes remained free throughout. 
Model~A starts with native widths ($w=1$) and free $\Delta t$ ($R^2 = 0.978$; Figure~\ref{fig:double_model_hierarchy}).
The high $R^2$ in Model~A is misleading: models A--D do not recover the two-peak envelope.
The characteristic trough appears only when the two normal components are allowed to widen.
Model~E reaches $R^2 = 0.999$ with free amplitudes ($a_1 = 0.88$, $a_2 = 0.94$), free $\Delta t = 1.60$\,ms, and a shared width scale $w = 1.29$.
Relative to the median normal half width of $1.58$\,ms, that scale corresponds to fitted component half widths of $2.04$\,ms each---wider than a native normal, not narrower.
Model~F, with independent width scales, fits equally well ($R^2 = 1.000$; component half widths $1.89$\,ms and $2.21$\,ms) but does not change the picture, so E is the working description.
The double is therefore not just a normal pulse repeated about 1\,ms later: recovering the trough requires broader components than the normal prototype.

\begin{figure}[H]
  \centering
  \includegraphics[width=0.85\textwidth]{pulse_shapes/normal_pulses_aligned_to_double.png}
  \caption{Two median normal pulses aligned to the peaks of the median double pulse. 
  The solid red trace is the median double, with red stars on its peaks. 
  The dashed and dotted blue curves are the median normal, shifted along the x-axis so each sits under one of those peaks. 
  The trough lands where the two normals cross.
  A pure short doublet of unmodified normals already looks similar, but the quantitative fits (Figure~\ref{fig:double_model_hierarchy}) show that matching the trough requires broader components.}
  \label{fig:normal_aligned_double}
\end{figure}

\begin{figure}[H]
  \centering
  \includegraphics[width=0.95\textwidth]{pulse_shapes/double_pulse_model_hierarchy_median.png}
  \caption{Six ways to model the median double pulse from two median normal components (constraints in Table~\ref{tab:double_pulse_models}).
  Each panel is labelled by model letter and $R^2$ and shows the median double (red, stars on the peaks), the model sum (grey), and the two normal components C1 and C2 (dashed and dotted blue).
  From A to F the constraints on $\Delta t$ and $w$ get looser ($w=1$ is the native normal width; $w>1$ widens).
  A already reaches $R^2 = 0.978$ with native widths.
  E and F let widths float and push $R^2$ near 1, with fitted $\Delta t$ drifting away from the observed 1\,ms peak spacing and with components broader than the native normal.
  The characteristic trough of the double pulse shape is recovered only by models E and F.}
  \label{fig:double_model_hierarchy}
\end{figure}

\subsection{Activity across timescales}

Pulse rates normalized for recording effort are shown separately for all EODs, wide pulses, and double pulses at four timescales in Figure~\ref{fig:activity_by_shape}.
Across the full dataset the effort-normalized mean rates are $1.21$\,Hz across all EODs.
The special pulse shape rates are $0.38$\,Hz for wide, and $0.018$\,Hz for double pulses.
Medians in the panels are active-session medians with bootstrap 95\% confidence interval (CI) shading: only sessions with a positive pulse rate in that bin enter the median, so empty or silent sessions are excluded.

On the 24\,h scale (Figure~\ref{fig:activity_by_shape}, top row), we can observe th circadian rhythm. 
The all pulse rate sits near $\sim$0.6\,Hz through the night (00:00 to 04:00), rises after about 05:00, and peaks near 08:00 to 09:00 at roughly 1.55\,Hz. 
After that morning peak the median falls to about 1.1\,Hz by 10:00 and ranges between roughly 0.75 and 1.0\,Hz for the rest of the day. 
Double and wide medians stay near zero across the clock, with wide pulses showing the same pulse rate elevation in the morning as the all pulse curve. 
That morning maximum lines up with feeding related activity (Section~\ref{sec:feeding_corr}).

The calendar month scale follows a similar pattern (Figure~\ref{fig:activity_by_shape}, second row).
All pulse medians sit around 0.6 to 1.0\,Hz for most of the year and there is one clear peak in November ($\sim$1.7\,Hz) that decreases in December and returns to baseline in January. 
There are two noticable peaks in the confidence band at March and May rate increase in March.
As seen in the 24h scale, wide pulses track that November bump from all pulses. 
They are near zero from February through September, rise in October, and reach about 0.6\,Hz in November before returning to baseline in January.
Double pulses stay at a near-zero baseline of roughly $0.018$\,Hz all year without any noticeable peaks or patterns.

Months since recording start (Figure~\ref{fig:activity_by_shape}, third row) adds more information to explain the increased late year pulse rates in the calendar month view.
There is one peak in November 2025 ($\sim$2.4\,Hz), indicating that the November bump is not seasonal but can be attributed to specific behavior in 11/2025. 
Beyond that, early 2024 carried the highest all pulse medians, with peaks near $\sim$6\,Hz in April 2024 and $\sim$5.6\,Hz in June 2024.
It is important to add here that the CIs are wide in those months. 
Also, these peaks are not reflected in the calendar month view although the November 2025 peak is. 
It is likely that they are driven by a few sessions with high pulse rates rather than a representative increase in activity across sessions.
After June 2024 the all pulse median falls sharply: July 2024 is a single overnight session near $\sim$0.9\,Hz, and August to October 2024 remain near $\sim$0.7 to 1.8\,Hz.
Gaps appear where no sessions contribute. 
This is the case for January 2024 and February 2025.
Wide pulses are almost absent until the November 2025 spike ($\sim$1.4\,Hz). 
Double pulses stay low, with a small rise towards the end of the recordings on the order of 0.5\,Hz into early 2026. 

Yearly bins (Figure~\ref{fig:activity_by_shape}, bottom row) show a steady drop in all pulse median rate, from about 3.5\,Hz in 2023 through 1.7\,Hz (2024) and 1.0\,Hz (2025) to about 0.5\,Hz in 2026.
The bootstrap band also narrows over those years.
Dual-line weighting (Methods) is already applied in these session medians, so the drop is not an uncorrected double count from the second electrode line.
2023 and 2026 are still partial calendar years, so those two points sample only the months that were recorded.
Double and wide yearly medians remain near $\sim$0.1\,Hz and do not follow that long-term decline. 

\begin{figure}[H]
  \centering
  \includegraphics[width=\textwidth]{activity_timescales/circadian_panels_by_shape_stacked.png}
  \caption{Pulse rate across different timescales, shown separately for all pulses (left), wide pulses (middle), and double pulses (right).
  Thick traces are active session medians; washed-out bands show bootstrap 95\% CIs.
  Rows from top to bottom: 24\,h (hourly bins), calendar months, months since recording start, and yearly bins from 2023 to 2026.}
  \label{fig:activity_by_shape}
\end{figure}

\subsection{Environmental and feeding correlations}
\label{sec:feeding_corr}

Monthly means of pulse rate against tank temperature and conductivity show little structure in the measured ranges (Figure~\ref{fig:env_corr}; $n=13$ months with overlapping sensor logs). 
Temperature spans only about 25.4 to 26.8\,$^{\circ}$C. 
Conductivity roughly spans from 180 to 350\,$\mu$S\,cm$^{-1}$. 
Pearson tests on the plotted linear fits are flat for all pulses (temperature: $r=-0.06$, $p=0.84$; conductivity: $r=0.01$, $p=0.96$) and for wide pulses (temperature: $r=-0.14$, $p=0.65$; conductivity: $r=-0.38$, $p=0.20$). 
Double-pulse rate vs temperature is the only significant linear association ($r=0.62$, $p=0.023$), but the matching Spearman rank correlation is not ($\rho=-0.07$, $p=0.81$). 
So that correlation is not robust to rank order and likely reflects a few leverage points in a small sample. 
Double vs conductivity is likewise non-significant ($r=0.23$, $p=0.44$). 
Wide and double rates sit much lower than the all-pulse scatter.
With only $n=13$ months, the informative value of this analysis is limited meaning that existing correlations could be missed. 

\begin{figure}[H]
  \centering
  \includegraphics[width=\textwidth]{correlations/correlations_all_pulse_shapes.png}
  \caption{Environmental correlations of pulse rate in Hz with tank temperature in $^{\circ}$C (left) and conductivity in µS\,cm$^{-1}$ (right). 
  Each cloud is a pulse category (purple - all pulses, red - double, green - wide), with a dashed line of the same colour indicating the linear fit ($n=13$ months). 
  The only significant Pearson correlation is double vs temperature ($p=0.023$).
  In a Spearman rank correlation all three traces are not significant.
  Special-shape rates sit well below the all-pulse scatter.}
  \label{fig:env_corr}
\end{figure}

Correlations of activity and feeding times show a much clearer picture. 
Feeding windows were defined as $\pm$5\,min around the logged feeding time, and baseline windows as $>$5\,min away from feeding.
Mean pulse rate for all pulses is $1.34$\,Hz during non-feeding times, versus a $3.8$-times increase to $5.03$\,Hz in feeding windows (Figure~\ref{fig:feeding_bars}; two-sided Mann--Whitney $U$, $p<0.001$). 
Wide pulses rise from $0.45$\,Hz to $0.64$\,Hz ($1.4$-times; $p<0.001$). 
Double-pulse window means stay near $0.03$\,Hz ($0.033$\,Hz baseline vs $0.030$\,Hz in feeding windows), so the mean contrast does not show a net increase.
Still, the minute-level distributions differ significantly($p<0.001$) because doubles occupy a larger fraction of feeding minutes.
On its own peri-event scale (Figure~\ref{fig:feeding_traj}), the double-pulse rate shows a brief peak near $0.05$\,Hz at the logged feeding time from a near-zero pre-feeding baseline.
Peri feeding trajectories for all pulses show the rise starting about 2\,min before the logged feeding time, peaking near $9$\,Hz at time zero, then decaying over the next 10\,min.
Wide pulses follow the same shape at lower amplitude (peak near $1.1$\,Hz).
Mean feeding clock time is 08:34 (bootstrap 95\% CI 08:14--08:57; $n=127$ sessions). 
That morning rate peak in Figure~\ref{fig:activity_by_shape} coincides with feeding times.

\begin{figure}[H]
  \centering
  \includegraphics[width=0.75\textwidth]{correlations/feeding_vs_nonfeeding_pulse_rates.png}
  \caption{Mean pulse rate in feeding windows ($\pm$5\,min) next to the non feeding baseline ($>$5\,min away). 
  Purple, green, and red are all, wide, and double pulses. 
  Within each hue the darker bar is feeding and the lighter one is baseline. 
  Brackets mark two-sided Mann--Whitney $U$ tests on minute-level rates ($***$: $p<0.001$). 
  Overall activity and wide pulses increase during feeding; doubles stay near the x-axis on this shared mean scale but still differ by the rank test.}
  \label{fig:feeding_bars}
\end{figure}

\begin{figure}[H]
  \centering
  \includegraphics[width=0.85\textwidth]{correlations/peri_feeding_pulse_rate_trajectories.png}
  \caption{Minute by minute rate around feeding, one panel each for all pulses, wide pulses, and double pulses, each with its own y-axis.
  Coloured traces are means across matched feeding events (purple - all pulses, green - wide, red - double).
  The matching washed out band is $\pm$ SEM over those events (\texttt{scipy.stats.sem} in \texttt{correlate\_activity\_with\_feeding.py}). 
  The dashed vertical line is the logged feeding time that each trace was aligned to. 
  Overall rate starts climbing about two minutes early, hits roughly $9$\,Hz at time zero, then decays until the end of the time window at +10 minutes. 
  Wides follow the same shape at lower amplitude (peak near $1.1$\,Hz). 
  On the double-pulse panel a short peak near $0.05$\,Hz appears at time zero.}
  \label{fig:feeding_traj}
\end{figure}

\subsection{Position estimation}

Head position was taken as the electrode with the largest positive peak in each multi channel snippet. 
Collapsed over all hourly bins, occupancy is strongly end biased (Figure~\ref{fig:pos_overall}). 
Pulse rate peaks at electrode 0 (bright end, $\sim$0.33\,Hz) and electrode 15 (dark end, $\sim$0.36\,Hz). 
Mid bright electrodes (1 to 9) stay low. 
Dark side electrodes 10 to 14 carry more mid line activity than the bright mid line. 
The animals park at tank ends and use the dark half more while moving between ends.

\begin{figure}[H]
  \centering
  \includegraphics[width=0.85\textwidth]{position_estimation/overall_position_distribution_peak.png}
  \caption{Where along the 16 electrode line the pulses land when hourly bins are collapsed. 
  Yellow bars give pulse rate per electrode. 
  The dashed vertical line is the bright/dark cut at 2.5\,m, and the light grey wash marks the dark half (electrodes 10 to 15). 
  Animals spend most time at the two ends (0 and 15), and the dark mid electrodes stay busier than the bright mid stretch.}
  \label{fig:pos_overall}
\end{figure}

Session wise median head position across timescales is shown in Figure~\ref{fig:pos_panels}. 
On the 24\,h scale the median stays mostly on the bright side of the 2.5\,m boundary, dipping toward the front (lower meters) around 07:00 to 08:00 and sitting closer to the dark boundary at other hours. 
The median alone understates how often the upper percentile band crosses into the dark half. 
Calendar month and months since start show slow drifts of preferred position, including a late 2025 push toward the dark boundary.

\begin{figure}[H]
  \centering
  \includegraphics[width=\textwidth]{position_estimation/position_panels_peak.png}
  \caption{Head position through the same four timescales as the activity figure. 
  The yellow line is the session wise median, the light yellow band the 16th to 84th percentile, and the faint yellow dots the session values themselves. 
  Everything above the dashed line (grey wash) is the dark half of the tank past 2.5\,m. 
  Medians usually sit on the bright side, slide toward the front mid morning, and wander over months, including a late 2025 stretch nearer the dark boundary.}
  \label{fig:pos_panels}
\end{figure}

Figure~\ref{fig:pos_anim} shows selected frames from a short turnaround scene in one recording.
An eel icon is plotted to move along the electrode line with the head position equal to the strongest positive electrode.
The printed figure is a two-panel keyframe strip; the full short animation (GIF) is available via the QR code.

\begin{figure}[H]
  \centering
  \includegraphics[width=\textwidth]{position_estimation/eel_position_keyframes.png}

  \vspace{0.6em}
  \begin{minipage}[c]{0.14\textwidth}
    \centering
    
\includegraphics[width=\linewidth]{position_estimation/eel_position_animation_qr.png}
  \end{minipage}%
  \hfill
  \begin{minipage}[c]{0.84\textwidth}
    \raggedright\footnotesize
    Scan the QR code for the head-position animation (GIF).
    Direct GIF:
    \url{https://raw.githubusercontent.com/saraheisele/mscthesis/master/docs/latex_thesis/figures/position_estimation/eel_position_animation.gif}.
  \end{minipage}

  \caption{Two frames from one Berlin recording (dark-area logger chunk
  \texttt{eellogger02-20260128T163148.wav}).
  Panel~(a): head at electrode~0, facing the bright end; the raw trace shows
  electrode~0 with a large positive EOD.
  Panel~(b): head at electrode~15, facing the dark end; the same electrode~0
  trace is polarity-flipped and much smaller.
  Head placement uses the electrode with the strongest positive peak.
  In this printed figure both raw panels deliberately share fixed electrode~0
  so amplitude and polarity can be compared on one channel; in the GIF
  animation the right-hand raw panel instead follows each pulse's head
  electrode (the channel with the strongest positive peak at that time).
  The bright (yellow) and dark (grey) tank halves are indicated by the overlapping circles,
  with the dashed line at the bright/dark boundary.
  The QR code opens the matching GIF supplement.}
  \label{fig:pos_anim}
\end{figure}

\newpage
